\section{Methodology}
\label{sec:method}

Given a scene layout generated by a 3D scene generation baseline (e.g., DiffuScene ~\cite{tang2024diffuscenedenoisingdiffusionmodels}), represented as a set of oriented bounding boxes, our goal is to produce a collision-free 3D scene that satisfies fundamental physical constraints.

To this end, we introduce a fine-grained collision resolution pipeline that eliminates object interpenetrations through alternating room-bound containment and Separating Axis Theorem (SAT)-based collision repair, as illustrated in 
Figure~\ref{fig:pipeline}. Given a text prompt describing the desired scene, the baseline generation model first synthesizes an initial layout of oriented bounding boxes and retrieves or generates the corresponding 3D object meshes. These 3D object meshes are positioned, scaled, and rotated according to the predicted layout. Our \emph{Collision Resolver} then operates purely on this set of oriented bounding boxes, iteratively alternating between room-bound containment and SAT-based collision repair until convergence, i.e., until no collisions between objects remain. The refined bounding-boxes are finally reapplied to the retrieved meshes, yielding a physically plausible, collision-free 3D scene that preserves the semantic layout intended by the baseline model.
\label{sec:method:prelim}
\begin{figure*}[t]
    \centering
    \includesvg[width=\textwidth,inkscapelatex=false]{Figure1_ACOMPA.svg}
    \caption{Overview of the proposed Collision Resolver. Given a user-provided text prompt, a 3D scene generation baseline (eg., DiffuScene)will first generates a scene layout represented by oriented bounding boxes and retrieves the corresponding 3D object meshes. Our collision resolver then iteratively alternates between roombound containment and OBB collision repair using the Separating Axis Theorem (SAT) until convergence, i.e., no collisions remain in the scene. The refined bounding boxes are finally used to produce a physically plausible, collision-free 3D scene while preserving the original semantic layout. The bottom-left panel shows an input layout generated by DiffuScene, where the red boxes highlight object interpenetrations and out-of-bound placements. The bottom-right panel shows the corresponding layout after applying our SAT-based collision resolver, with the resolved regions highlighted in green.}
    \label{fig:pipeline}
\end{figure*}

\subsection{Preliminaries}
DiffuScene~\cite{tang2024diffuscenedenoisingdiffusionmodels} formulates indoor scene synthesis as a conditional denoising diffusion process over scene layouts. Starting from Gaussian noise, the diffusion model progressively predicts a clean set of object parameters conditioned on the room type and scene context. At the end of the reverse diffusion process, the model produces a structured scene layout consisting of $M$ oriented bounding boxes, each corresponding to a furniture instance.

Each bounding box is parameterized as
\begin{equation}
\mathbf{b}_i=(l_i,h_i,w_i,x_i,y_i,z_i,\sin\theta_i,\cos\theta_i)\in\mathbb{R}^{8},
\end{equation}
where $(x_i,y_i,z_i)$ denotes the center position $\mathbf{p}_i$, $(l_i,h_i,w_i)$ denotes the object dimensions, and $\theta_i$ is the yaw orientation. We further define an \emph{objectness mask} $\mathbf{o}\in\{0,1\}^{M}$ to distinguish valid furniture instances ($o_i=1$) from auxiliary nodes such as the floor and room boundary.

\subsection{Post-Processing Collision Resolver}
\label{sec:method:resolver}

Because diffusion-based layout generators such as DiffuScene do not enforce any implicit or explicit physical constraints during the denoising process, the resulting bounding-box layouts frequently exhibit inter-object penetrations and out-of-bound placements. We address this with our \emph{Collision Resolver} - a deterministic post-processing stage (Algorithm~\ref{alg:collision_resolver}) 
that iteratively removes inter-object collisions and invalid object placements, producing a physically-plausible 3D scene directly at inference time.

The resolver first performs a fast 1D vertical bounding-box check along the $y$-axis: object pairs that are vertically separated (e.g., a monitor resting on a desk) require no further correction and are immediately excluded. For 
every remaining objects, vertically overlapping pair, resolution proceeds entirely in the 2D floor plane ($xz$), combining (a) \emph{oriented room-bound containment clamping}, and (b) \emph{OBB-SAT collision detection and resolution}, which generalizes the classical Minimum Translation Vector (MTV) through an alignment-biased axis selection.

\begin{algorithm}[!t]
\caption{Post-Processing SAT Collision Resolver}
\label{alg:collision_resolver}                        
\begin{algorithmic}[1]
\REQUIRE Layout boxes $\{\mathbf{b}_i\}_{i=1}^{M}$ with centers $\mathbf{p}_i$; mask $\mathbf{o}$; class indices $\{\ell_i\}$; max iterations $K_{\max}=500$; clearance $\varepsilon_{\mathrm{push}}$; bias $\beta$
\ENSURE Collision-free layout $\{\mathbf{b}_i'\}$
\STATE $\mathcal{L} \gets \{i : o_i = 0 \ \lor\ \ell_i = 0\}$ \COMMENT{auxiliary floor/layout nodes}
\STATE $L^\ast \gets \operatorname*{arg\,max}_{i \in \mathcal{L}} \; l_i \, w_i$ \COMMENT{largest-footprint reference room box}
\FOR{$k = 1, 2, \ldots, K_{\max}$}
    \STATE $\textit{converged} \gets \textbf{true}$
    \STATE \COMMENT{1. Oriented room-bound containment (Section~\ref{sec:method:room_containment})}
    \IF{$\mathcal{L} \neq \varnothing$}
        \FOR{each object $i \notin \mathcal{L}$}
            \STATE Project corners $\mathcal{C}_i$ into the room frame of $L^\ast$ via Eq.~\eqref{eq:room_proj}
            \STATE Compute correction $(\Delta u, \Delta v)$ via Eq.~\eqref{eq:room_clamp}
            \IF{$(\Delta u, \Delta v) \neq (0,0)$}
                \STATE Update $\mathbf{p}_i$ via Eq.~\eqref{eq:room_update}; \ $\textit{converged} \gets \textbf{false}$
            \ENDIF
        \ENDFOR
    \ENDIF
    \STATE \COMMENT{2. Object-object collision resolution}
    \FOR{each pair $(i, j)$ with $i < j$, $o_i = o_j = 1$}
        \IF{vertically separated (1D $y$-axis check)}
            \STATE \textbf{continue} 
        \ENDIF
        \STATE Compute corner sets $\mathcal{C}_i, \mathcal{C}_j$ in $xz$-plane
        \STATE Compute overlap depths $\{\delta_k\}_{k=1}^{4}$ via Eq.~\eqref{eq:sat_overlap}
        \IF{$\exists\, k: \delta_k \leq 0$}
            \STATE \textbf{continue} \COMMENT{no collision}
        \ENDIF
        \STATE $\hat{\mathbf{c}}_{ij} \gets (\mathbf{p}_j - \mathbf{p}_i) / \|\mathbf{p}_j - \mathbf{p}_i\|$
        \STATE Select push axis $\hat{\mathbf{a}}^*$ and depth $\delta^*$ via Eq.~\eqref{eq:axis_score}
        \STATE $\Delta \gets (\delta^* + \varepsilon_{\mathrm{push}}) / 2$
        \STATE $\mathbf{p}_i \gets \mathbf{p}_i - \Delta \cdot \hat{\mathbf{a}}^*$
        \STATE $\mathbf{p}_j \gets \mathbf{p}_j + \Delta \cdot \hat{\mathbf{a}}^*$     
        \STATE $\textit{converged} \gets \textbf{false}$
    \ENDFOR
    \IF{$\textit{converged}$}
        \STATE \textbf{break}
    \ENDIF
\ENDFOR
\RETURN $\{\mathbf{b}_i'\}$
\end{algorithmic}
\end{algorithm}
\subsubsection{Oriented Room-Bound Containment}
\label{sec:method:room_containment}

Object placements produced by 3D diffusion models are not guaranteed to lie within the room footprint, and a naive axis-aligned clamp is incorrect whenever the room itself is rotated with respect to the world frame. We therefore constrain every furniture instance within the room boundary expressed in the room's \emph{own} oriented frame, rather than in world coordinates.

We first isolate the auxiliary layout nodes $\mathcal{L} = \{i : o_i = 0 \ \lor\ \ell_i = 0\}$, where $\ell_i = \arg\max_c \mathbf{y}_i^{(c)}$ is the predicted class index of box $i$ and class $0$ denotes floor/room instances. When a scene contains several such instances, we take the one with the largest footprint area,
\begin{equation}
L^{\ast} = \operatorname*{arg\,max}_{i \in \mathcal{L}} \; l_i \, w_i,
\label{eq:room_select}
\end{equation}
and read off its center $\mathbf{c}_L$, half-extents $(h_{xL}, h_{zL}) = (l_{L^{\ast}}/2, w_{L^{\ast}}/2)$, and oriented axes $\hat{\mathbf{a}}_L = (\cos\theta_{L^{\ast}}, \sin\theta_{L^{\ast}})$, $\hat{\mathbf{a}}_L^{\perp} = (-\sin\theta_{L^{\ast}}, \cos\theta_{L^{\ast}})$. Alternatively, for models such as DiffuScene that are conditioned on a fixed room geometry, we directly compute the bounding box of the ground-truth floor plan vertices, yielding an axis-aligned room frame with center $\mathbf{c}_L = (0, 0)$, half-extents $(h_{xL}, h_{zL})$, and identity axes $\hat{\mathbf{a}}_L = (1, 0), \hat{\mathbf{a}}_L^{\perp} = (0, 1)$. Together, these define the room's local frame. If neither a layout node nor a floor plan is available, the containment step is skipped for that scene.

For every non-layout object $i \notin \mathcal{L}$, we express its four floor-plane corners $\mathcal{C}_i$ in this local frame,
\begin{equation}
u = (\mathcal{C}_i - \mathbf{c}_L) \cdot \hat{\mathbf{a}}_L, \qquad
v = (\mathcal{C}_i - \mathbf{c}_L) \cdot \hat{\mathbf{a}}_L^{\perp},
\label{eq:room_proj}
\end{equation}
and compute a per-axis correction $(\Delta u, \Delta v)$ that pulls any out-of-bound corners back within the half-extents. Along $u$,
\begin{equation}
\Delta u =
\begin{cases}
-\dfrac{u_{\max} + u_{\min}}{2}, & u_{\max} - u_{\min} > 2 h_{xL}, \\[6pt]
h_{xL} - u_{\max}, & u_{\max} > h_{xL}, \\[3pt]
-h_{xL} - u_{\min}, & u_{\min} < -h_{xL}, \\[3pt]
0, & \text{otherwise},
\end{cases}
\label{eq:room_clamp}
\end{equation}
with $\Delta v$ defined analogously along $\hat{\mathbf{a}}_L^{\perp}$ using $h_{zL}$. The first branch handles objects whose footprint exceeds the room extent along that axis (e.g.\ an oversized asset), in which case the object is re-centered on the room rather than driven against a single wall; the remaining branches clamp the object to whichever wall it protrudes through, and leave objects that are already contained unchanged. The world-space center is then updated as
\begin{equation}
\mathbf{p}_i \gets \mathbf{p}_i + \Delta u\, \hat{\mathbf{a}}_L + \Delta v\, \hat{\mathbf{a}}_L^{\perp}.
\label{eq:room_update}
\end{equation}
This step is applied to all non-layout objects at the start of every outer iteration $k$ (Algorithm~\ref{alg:collision_resolver}), \emph{before} pairwise SAT collision resolution. As a result, the boundary-penalized overlap depth used by the alignment-biased axis selection (Section~\ref{sec:method:axis_selection}, Eq.~\eqref{eq:axis_score}) is always evaluated against an up-to-date, already-contained room configuration, so that object-object pushes and room clamping will cooperate rather than fight each other across iterations.

\subsubsection{OBB Collision Detection via SAT}
\label{sec:method:obb_sat}

We detect collisions between oriented bounding boxes (OBBs) using the Separating Axis Theorem (SAT) \cite{10.1145/237170.237244}. For two OBBs with yaw angles $\theta_i, \theta_j$, we test four candidate separating axes. 

Let $\mathbf{a}_k$ ($k = 1, \ldots, 4$) denote these axes, and let $\mathcal{C}_i, \mathcal{C}_j$ denote the sets of four floor-plane corners of boxes $i$ and $j$. For each axis $\mathbf{a}_k$, we project both sets of corners and compute the \emph{overlap depth}:
\begin{equation}
\begin{split}
\delta_k &= \min\!\big(\max(\mathcal{C}_i \cdot \mathbf{a}_k), \max(\mathcal{C}_j \cdot \mathbf{a}_k)\big) \\
&\quad - \max\!\big(\min(\mathcal{C}_i \cdot \mathbf{a}_k), \min(\mathcal{C}_j \cdot \mathbf{a}_k)\big).
\end{split}
\label{eq:sat_overlap}
\end{equation}
If $\delta_k \leq 0$ for any axis, the boxes are not colliding. Otherwise, all four axes yield positive overlaps.

\subsubsection{Alignment-Biased Axis Selection}
\label{sec:method:axis_selection}

The standard SAT approach selects the axis with the minimum overlap depth $\delta^* = \min_k \delta_k$. For interior layouts, this often produces unintuitive sliding. We introduce an \emph{alignment-biased} axis selection that prefers axes aligned with the center-to-center direction $\mathbf{c}_{ij} = \mathbf{p}_j - \mathbf{p}_i$. 

Each axis $\mathbf{a}_k$ is oriented to point from $i$ toward $j$: $\hat{\mathbf{a}}_k = \mathrm{sign}(\mathbf{a}_k \cdot \mathbf{c}_{ij}) \cdot \mathbf{a}_k$. We compute an alignment score:
\begin{equation}
\alpha_k = \big|\hat{\mathbf{a}}_k \cdot \hat{\mathbf{c}}_{ij}\big|, \qquad
s_k = \frac{\tilde{\delta}_k}{\alpha_k^{\beta} + \epsilon}
\label{eq:axis_score}
\end{equation}
where $\beta=2$ and $\epsilon$ is a small constant. Crucially, $\tilde{\delta}_k$ represents a \emph{boundary-penalized} overlap depth, computed with respect to the same reference room frame $L^\ast$ used in Section~\ref{sec:method:room_containment}. For an object with corners $\mathcal{C}$, define its total room-boundary violation as
\begin{equation}
g(\mathcal{C}) = \max(0, u_{\max}-h_{xL}) + \max(0, -h_{xL}-u_{\min}) + \max(0, v_{\max}-h_{zL}) + \max(0, -h_{zL}-v_{\min}),
\label{eq:oob}
\end{equation}
using the projections $(u,v)$ of Eq.~\eqref{eq:room_proj}. Before penalizing an axis, we tentatively apply the half-depth push $\pm(\delta_k/2)\hat{\mathbf{a}}_k$ to the corner sets $\mathcal{C}_i, \mathcal{C}_j$ and compare $g(\cdot)$ before and after: if either object's violation \emph{increases} ($g_{\text{after}} - g_{\text{before}} > 10^{-3}$), the axis is deemed to push an object further out of the room, and we set $\tilde{\delta}_k = 1000\,\delta_k$; otherwise $\tilde{\delta}_k = \delta_k$. This preemptive check prevents the object-object resolution step from undoing the room containment of Section~\ref{sec:method:room_containment} at the next iteration. Axes with $\alpha_k < 0.15$ are discarded. The axis with the lowest score $s_k$ is selected as the push direction.



Given the selected push axis $\hat{\mathbf{a}}^*$ and depth $\delta^*$ for a colliding pair $(i, j)$, we symmetrically displace both objects by $(\delta^* + \varepsilon_{\mathrm{push}}) / 2$ along $\hat{\mathbf{a}}^*$:
\begin{equation}
\mathbf{p}_i \gets \mathbf{p}_i - \frac{\delta^* + \varepsilon_{\mathrm{push}}}{2}\, \hat{\mathbf{a}}^*, \qquad
\mathbf{p}_j \gets \mathbf{p}_j + \frac{\delta^* + \varepsilon_{\mathrm{push}}}{2}\, \hat{\mathbf{a}}^*
\label{eq:push}
\end{equation}
where $\varepsilon_{\mathrm{push}} = 0.02\,\text{m}$ provides clearance. Algorithm~\ref{alg:collision_resolver} iterates over all pairs until no collisions remain or $K_{\max} = 500$ is reached.